\section{Experiments Using U-TOE}
We conducted experiments to validate the functionality and compatibility of U-TOE on model side (universal support for model structure and ML frameworks) and on device side (wide-spectrum support for IoT devices). Hence, we dived into two orthogonal directions: For device support, we evaluated a quantized LetNet-5 on various IoT boards; For model compatibility, we evaluated multiple models on a local STM32F746G discovery board. 
\begin{itemize}
    \item \textbf{Model Selection} We selected pre-trained, quantized models from open source repositories\footnote{see \url{https://github.com/ARM-software/ML-zoo} and \url{https://github.com/mlcommons/tiny}}, which target on typical TinyML tasks (Image Classification, Key Word Spotting, Visual Word Wake, Noise Suppression and Abnormal Detection). The weights and activations of the model were quantized to 8-bit integer, yet the inputs and outputs remain in IEEE FP32 format.
    \item \textbf{Model Optimization} We only used built-in, rule-based optimization in uTVM. Thus, all heuristic optimization strategies like model scheduling were disabled.
    \item \textbf{MCU Configuration} We disabled data and instruction cache to observe the "memory wall" effect in ML model. The core clock frequency was pre-set by CPU initialization code in RIOT and is presented with experiment results in Section ~\ref{sec:results}. 
    \item \textbf{Hybrid Deployment} The experiments were conducted both on local and remote IoT boards provided by FIT-IoT Lab. 
\end{itemize}
It is noted that for each evaluation we preset the number of trials to ten in order to address random error.

\section{Analysing U-TOE Measurements}
\label{sec:results}
\subsection{Per-Model Evaluation}
Table \ref{tab:per-model-eval-result-multiboard} presents the resource consumption of LeNet-5 model on various IoT boards, generated by Per-Model evaluation. ARM Cortex-M series MCUs show no significant difference on memory and storage usage, and the computational latency declines as the core frequency increases. Benefits from fully support of DSP and Thumb-2 instruction set, Cortex-M3, -M4 MCUs perform better than Cortex-M0+ with the same core clock frequency. An \textbf{outlier} was discovered on SiFive RISC-V MCU. With the highest core clock frequency this won the least favorable ranking on computational latency and the memory usage is higher than other ARM Cortex-M MCU. The SiFive RISC-V MCU uses an external, SPI NOR flash for data and program storage, causing a huge performance regression while we disabled the cache.

Table \ref{tab:per-model-eval-result-multimodel} presents the results of various ML models on representative TinyML tasks on single IoT board, proving the universal support of ML framework and model structure. Except for LeNet-5 trained on local host device with Pytorch, all the others came from open source model zoo.
\begin{table*}[ht]
\centering
\begin{threeparttable}
\caption{Evaluation of LeNet-5 model on various IoT boards.}
\label{tab:per-model-eval-result-multiboard}
\centering
% TODO: The order of Core and Frequencies
\begin{tabular}{llllllll}
\hline
\multirow{2}{*}{Board   / MCU} & \multirow{2}{*}{Core} & \multirow{2}{*}{Memory (KB)} & \multirow{2}{*}{Storage (KB)} & \multicolumn{4}{l}{Computational Latency   (ms)}    \\ \cline{5-8} 
                               &                       &                              &                               & 95\%-CI               & Median  & Min.    & Max. \\\hline
b-l072z-lrwan1 / STM32L072CZ & M0+ @ 32 MHz          & 11.288                       & 64.34                         & {[}261.829, 262.249{]} & 262.187 & 261.35  & 262.216 \\
samr21-xpro / ATSAMR21G18A                   & M0+ @ 48 MHz          & 11.292                       & 64.956                        & {[}182.058, 182.083{]} & 182.068 & 182.04  & 182.097 \\
samr30-xpro / ATSAMR30G18A & M0+ @ 48 MHz          & 11.208                       & 65.168                        & {[}176.936, 176.965{]} & 176.958 & 176.924 & 176.975 \\
samr34-xpro / ATSAMR34J18   & M0+ @ 48 MHz          & 11.296                       & 65.436                        & {[}178.686, 178.718{]} & 178.708 & 178.669 & 178.732 \\
arduino-zero   / ATSAMD21G18        & M0+ @ 48 MHz          & 11.292                       & 64.94                         & {[}182.061, 182.082{]} & 182.068 & 182.051 & 182.098 \\
rpi-pico   / RP2040        & M0+ @ 125 MHz   &   28.704     &    109.504  &  [70.108, 70.13]  &   70.117  &   70.091 & 70.151 \\

openmote-b   / CC2538SF53          & M3 @ 32 MHz           & 11.1                         & 66.08                         & {[}200.337, 200.384{]} & 200.367 & 200.323 & 200.404 \\
IoT-LAB M3 / STM32F103REY    & M3 @ 72 MHz           & 11.296                       & 62.26                         & {[}97.74, 97.757{]}    & 97.751  & 97.733  & 97.764  \\

nucleo-wl55jc   / STM32WL55JC  & M4 @ 48 MHz           & 11.288                       & 63.18                         & {[}98.649, 98.668{]}   & 98.661  & 98.637  & 98.679  \\
nrf52dk   / nRF52832           & M4 @ 64 MHz           & 11.328                       & 61.012                        & {[}66.124, 66.152{]}   & 66.132  & 66.096  & 66.158  \\
nrf52840dk   / nRF52840        & M4 @ 64 MHz           & 11.348                       & 61.332                        & {[}66.078, 66.112{]}   & 66.088  & 66.087  & 66.163  \\
b-l475e-iot01a   / STM32L475VG & M4 @ 80 MHz           & 11.288                       & 61.604                        & {[}52.9, 52.901{]}     & 52.901  & 52.9    & 52.902  \\
% dwm1001   / nRF52832           & M4 @ 64 MHz           & 11.328                       & 61.012                        &                       &         &         &         \\
stm32f746g-disco / STM32F746NG       & M7 @ 216 MHz          & 11.076                       & 64.712                        & {[}39.6, 39.602{]}     & 39.601  & 39.599  & 39.604 \\\hline

esp32-wroom-32 / ESP32-D0WDQ6  & ESP32 @ 80 MHz       & 115.958                       & 157.719                        & {[}85.58, 85.583{]} & 85.582  & 85.576 & 85.584
\\\hline
esp32c3-devkit / ESP32-C3FN4 &  RISC-V @ 80 MHz      &    258.874       &    222.272   & [54.947, 54.957]   &      54.953    &   54.938   &    54.961 \\
sipeed-longan-nano / GD32VF103CBT6  & RISC-V @ 108 MHz &     103.108     &    106.422 & [37.783, 37.789]   &   37.789   &    37.779    &   37.791 \\
hifive1b / SiFive FE310-G002  & RISC-V @ 320 MHz       & 60.884                       & 66.492                        & {[}153.621, 154.166{]} & 153.747 & 153.717 & 154.938 \\

\hline
\end{tabular}
\begin{tablenotes}
    \item "Memory" and "Storage" columns refer to the resource consumption.
\end{tablenotes}
\end{threeparttable}
\end{table*}

\begin{table*}[ht]
\centering
\begin{threeparttable}
\caption{Evaluation of various quantized models on stm32f746-disco board.}
\label{tab:per-model-eval-result-multimodel}
\begin{tabular}{llllllll}
\hline
\multirow{2}{*}{Model} & \multirow{2}{*}{Task} & \multirow{2}{*}{Memory (KB)} & \multirow{2}{*}{Storage (KB)} & \multicolumn{4}{l}{Computational Latency   (ms)} \\ \cline{5-8}
                       &                       &                              &                               & 95\%-CI      & Median      & Min.     & Max.     \\\hline
DS-CNN   Small INT8\tnote{**}    & Keyword Spotting      &    68.992       &    71.796         &   [461.395, 461.396]  &    461.396         &   461.396       &   461.397       \\
MobileNetV1-0.25x INT8\tnote{*}         & Visual Wake Words     &  185.352      &    491.668         &  [1435.937, 1435.938]  &   1435.938           &  1435.938         &    1435.939      \\
LeNet-5 INT8          & Image Classification  &     12.068      &    65.851             &    [39.599, 39.603]          &   39.601   &   39.598       &   39.605     \\
Deep AutoEncoder INT8\tnote{*}         & Anomaly Detection     &    6.532     &   292.696 & [35.637, 35.638]    &     35.638   &    35.638   &    35.639   \\
RNNoise INT8\tnote{**}           & Noise Suppression     &   4.688       &    119.652        &   [12.151, 12.157]    &   12.154   &   12.148       &   12.160      \\\hline
\end{tabular}
\begin{tablenotes}
\item [*] These models originate from MLPerf Tiny Benchmarks repository on \url{https://github.com/mlcommons/tiny}.
\item [**] These models originate from ARM Model Zoo on \url{https://github.com/ARM-software/ML-zoo}.
\item All models were pre-trained and quantized by TFLite, except LeNet-5 INT8 was by Pytorch.
\item "Memory" and "Storage" columns refer to the resource consumption.
\end{tablenotes}
\end{threeparttable}
\end{table*}

\subsection{Per-Operator Evaluation}
We here used a tiny model with only three layers from TFlite as an example to avoid unnecessary complexity in demonstration, with output results presented in Table \ref{tab:per-ops-eval-result}. The computational bottlenecks are located in operator \emph{nn\_dense\_add\_nn\_relu} and \emph{nn\_dense\_add\_nn\_relu\_1}, and with the highest memory and storage consumption as well. We can trace down the corresponding layers of original model with the hints of associated parameters, which are the weights, bias or other trainable parameters of the model, making it possible to apply optimization strategies on well-targeted layers. 
\begin{table*}[]
\centering
\begin{threeparttable}
\caption{Per-Operator Evaluation Output of TFlite sinus model on stm32f746-disco board.}
\label{tab:per-ops-eval-result}
\begin{tabular}{llllll}
\hline
Ops                                                 & Time (us) & Time (\%) & Asso. Params                                        & Memory (KB) & Storage (KB) \\\hline
fused\_nn\_dense\_add\_nn\_relu    & 8.856        & 15.22\%   & p0, p1                 & 0.128       & 0.128        \\
fused\_nn\_dense\_add\_nn\_relu\_1 & 46.682        & 80.23\%   & p2, p3                             & 0.128       & 1.088        \\
fused\_nn\_dense\_add              & 2.646         & 4.54\%   & p4, p5 & 0.068       & 0.068       \\\hline
\end{tabular}
\begin{tablenotes}
    \item The uTVM auto-generated prefix of operator name \textit{tvmgen\_default\_} is not presented for the purpose of clarity.
\end{tablenotes}
\end{threeparttable}
\end{table*}


\section{Reproducible \& Custom U-TOE Experiments}
We released the full source code of the U-TOE toolkit on Github at \url{https://github.com/zhaolanhuang/U-TOE} under an open source LGPL v3 license. 

The implementations of all experiments can be found under \textit{examples} folders. Table \ref{tab:utoe-parameters} shows the parameters that can be tuned on user's purpose. For further details on how to start with U-TOE hands-on, the reader is referred to the comprehensive \textit{Readme.md} in the repository. 

On the one hand, researchers and practitioners who possess IoT hardware that is supported by the open source operating system RIOT (currently 250+ types of boards, using 60+ types of CPUs~\footnote{see \url{https://github.com/RIOT-OS/RIOT/tree/master/boards}} can use U-TOE out-of-the-box, directly on their boards. 

On the other hand, combined with the use of the free open-access testbed IoT-Lab~\footnote{see \url{https://www.iot-lab.info/}}, even researchers and practitioners who do not have such hardware on premises can conduct large scale experimental comparative evaluation campaigns using U-TOE.

{\bf Perspectives --} As RIOT board and CPU support expands and improves over time, and as uTVM also expands support to other architectures in parallel (both open source communities are very active) U-TOE can in a very short time expand its support for new use cases, automatically adding the support of uTVM for new boards, and the support of RIOT for new models. As such U-TOE may oragnically grow and become a useful link between the two communities.

Moreover, while the work on U-TOE in this paper has been focused on inference only on single-core microcontrollers, there is strong potential to extend the toolkit provided by U-TOE to support on-device learning scenarios, and for optimizing exploitation of multi-core microcontrollers.


\begin{table}[]
    \caption{Description of U-TOE Parameters.}
    \label{tab:utoe-parameters}
    \centering
    \begin{tabular}{cc}
    \hline
      Parameter   &  Remark \\\hline
      UTOE\_RANDOM\_SEED    & Random seed for randomized input \\
      UTOE\_TRIAL\_NUM & No. of trials in Per-Model evaluation \\
      UTOE\_RPC\_BUFFER\_SIZE & Buffer size of RPC server \\\hline 
      
    \end{tabular}
    
\end{table}